\documentclass[a4paper,12pt]{article}

\usepackage[T2A]{fontenc}
\usepackage[utf8]{inputenc}
\usepackage[russian]{babel}
\usepackage{amsmath}
\usepackage{graphicx}
\usepackage{booktabs}
\usepackage{geometry}
\usepackage{tabularx}
\geometry{margin=2cm}

\begin{document}
	
	\subsubsection*{Опыт 1. Образование слабого основания}
	
	\quad \textbf{ Ход опыта.}  В пробирку добавляют несколько капель раствора NH$_4$Cl, затем добавляют раствор NaOH с фенолфталеином.
	
	\textbf{Наблюдения.}   Розовая окраска исчезает; при подогревании ощущается запах аммиака.
	
	\textbf{Уравнения реакций:}
	\[
	NH_4Cl + NaOH \rightarrow NH_4OH + NaCl
	\]
	\[
	NH_4^+ + OH^- \rightleftharpoons NH_3\cdot H_2O
	\]
	
	\textbf{Расчёт константы равновесия.}
	
	Для системы:
	\[
	NH_4^+ + OH^- \rightleftharpoons NH_4OH
	\]
	
	константа равновесия равна:
	\[
	K = K_R(\text{NH}_4\text{OH}) = 10^{-5}.
	\]
	
	Значение $K = 10^{-5} \ll 1$, поэтому образование NH$_4$OH идёт лишь частично.  
	Равновесие сильно смещено в сторону исходных ионов, что подтверждает слабость основания и объясняет исчезновение окраски фенолфталеина.
	
	\textbf{Вывод.}  
	Образуется слабое основание NH$_4$OH; равновесие смещено влево, при нагревании выделяется NH$_3$.
	
	
	\subsubsection*{Опыт 2. Образование устойчивой слабой кислоты}
	
	\quad \textbf{ Ход опыта.}  
	К раствору H$_2$SO$_4$ добавляют ацетат натрия CH$_3$COONa.
	
	\textbf{Наблюдения.}  
	Появляется запах уксусной кислоты.
	
	\textbf{Уравнения:}
	\[
	H^+ + CH_3COO^- \rightleftharpoons CH_3COOH
	\]
	
	\textbf{Расчёт константы равновесия.}
	
	Для реакции:
	\[
	H^+ + CH_3COO^- \rightleftharpoons CH_3COOH
	\]
	
	используем константу диссоциации уксусной кислоты:
	\[
	K_R(\text{CH}_3\text{COOH}) = 10^{-5}.
	\]
	
	Тогда константа образования равна:
	\[
	K = \frac{1}{10^{-5}} = 10^5.
	\]
	
	\quad Поскольку $K \gg 1$, равновесие резко смещено в сторону образования CH$_3$COOH.  
	Это обосновывает появление запаха уксусной кислоты при подогреве.
	
	\textbf{Вывод.}  
	Образуется слабая кислота CH$_3$COOH; сдвиг равновесия обеспечивает наличие избытка ацетат-ионов.

	
	\subsubsection*{Опыт 3. Образование неустойчивой слабой кислоты}
	
	\quad \textbf{ Ход опыта.}  
	К раствору Na$_2$CO$_3$ добавляют раствор H$_2$SO$_4$.
	
	\textbf{Наблюдения.}  
	Интенсивное выделение CO$_2$.
	
	\textbf{Уравнения:}
	\[
	CO_3^{2-} + 2H^+ \rightarrow H_2CO_3 \rightleftharpoons CO_2\uparrow + H_2O
	\]
	
	\textbf{Расчёт константы равновесия.}
	
	Для реакции:
	\[
	CO_3^{2-} + 2H^+ \rightarrow H_2CO_3
	\]
	
	константа диссоциации нестойкой угольной кислоты:
	\[
	K_R(\text{H}_2\text{CO}_3) = 10^{-17}.
	\]
	
	Следовательно,
	\[
	K = \frac{1}{10^{-17}} = 10^{17}.
	\]
	  
	Очень большое значение $K$ означает, что H$_2$CO$_3$ неустойчива и тут же разлагается, переходя в CO$_2$ и воду.  
	Реакция протекает практически полностью, что соответствует интенсивному выделению газа.
	
	\textbf{Вывод.}  
	H$_2$CO$_3$ разлагается на CO$_2$ и воду, поэтому реакция идёт практически полностью.
	
	
	\subsubsection*{Опыт 4. Реакция нейтрализации}
	
	\quad\textbf{ Ход опыта.}  
	К раствору NaOH с фенолфталеином добавляют HCl.
	
	\textbf{Наблюдения.}  
	Исчезает розовая окраска.
	
	\textbf{Уравнение:}
	\[
	H^+ + OH^- \rightarrow H_2O
	\]
	
	\textbf{Расчёт константы равновесия.}
	
	Реакция:
	\[
	H^+ + OH^- \rightarrow H_2O
	\]
	
	её константа равна:
	\[
	K = \frac{1}{K_R(\text{H}_2O)} = \frac{1}{10^{-14}} = 10^{14}.
	\]
	
	\quad Поскольку $K$ огромно, реакция нейтрализации протекает практически полностью, что и приводит к исчезновению окраски фенолфталеина.
	
	\textbf{Вывод.}  
	Типичная реакция нейтрализации, протекает практически полностью.
	
	
	\subsubsection*{Опыт 6. Переосаждение малорастворимых веществ}
	
	\quad\textbf{ Ход опыта.}  
	Смешивают растворы Pb(NO$_3$)$_2$ и Na$_2$SO$_4$ (образуется PbSO$_4$). Затем добавляют K$_2$CrO$_4$.
	
	\textbf{Наблюдения.}  
	Белый осадок PbSO$_4$ превращается в жёлтый PbCrO$_4$.
	
	\textbf{Уравнения:}
	\[
	Pb^{2+} + SO_4^{2-} \rightarrow PbSO_4\downarrow
	\]
	\[
	Pb^{2+} + CrO_4^{2-} \rightarrow PbCrO_4\downarrow
	\]
	
	\textbf{Расчёт равновесия.}
	
	Произведения растворимости:
	\[
	\text{ПР}_{\text{PbSO}_4} = 10^{-8}, \qquad 
	\text{ПР}_{\text{PbCrO}_4} = 10^{-13}.
	\]
	
	Поскольку:
	\[
	10^{-13} \ll 10^{-8},
	\]
	
	то PbCrO$_4$ менее растворим, чем PbSO$_4$.
	
	Ион CrO$_4^{2-}$ вытесняет SO$_4^{2-}$ из осадка, приводя к переосаждению более малорастворимого жёлтого PbCrO$_4$.
	
	\textbf{Вывод.}  
	PbCrO$_4$ менее растворим, чем PbSO$_4$, поэтому происходит переосаждение.
	
	
	\subsubsection*{Опыт 7. Растворение осадков под действием сильной кислоты}
	
	\quad\textbf{ Ход опыта.}  
	Осадки FeS и CuS обрабатывают HCl.
	
	\textbf{Наблюдения.}  
	FeS растворяется с выделением H$_2$S, CuS практически не растворяется.
	
	\textbf{Уравнения:}
	\[
	FeS + 2H^+ \rightarrow Fe^{2+} + H_2S\uparrow
	\]
	\[
	CuS + 2H^+ \not\rightarrow 
	\]
	
	\textbf{Расчёт равновесий.}
	
	Произведения растворимости:
	\[
	\text{ПР}_{\text{FeS}} = 10^{-17}, \qquad
	\text{ПР}_{\text{CuS}} = 10^{-36}.
	\]
	
	Поскольку:
	\[
	10^{-36} \ll 10^{-17},
	\]
	
	CuS значительно менее растворим, чем FeS.
	 
	\quad FeS растворяется в кислоте, образуя H$_2$S.  
	CuS практически не растворяется из-за крайне малого ПР.
	
	\textbf{Вывод.}  
	CuS обладает значительно меньшей растворимостью; HCl не может его растворить.
	
	
	\subsubsection*{Опыт 8. Растворение осадка с образованием комплексного катиона}
	
	\quad\textbf{ Ход опыта.}  
	К раствору CuSO$_4$ добавляют NH$_4$OH: сначала выпадает Cu(OH)$_2$, затем при избытке NH$_3$ осадок растворяется.
	
	\textbf{Наблюдения.}  
	Осадок растворяется, раствор становится тёмно-синим.
	
	\textbf{Уравнения:}
	\[
	Cu^{2+} + 2OH^- \rightarrow Cu(OH)_2\downarrow
	\]
	\[
	Cu^{2+} + 4NH_3 \rightarrow [Cu(NH_3)_4]^{2+}
	\]
	
	\textbf{Расчёт равновесия.}
	
	Произведение растворимости:
	\[
	\text{ПР}_{\text{Cu(OH)}_2} = 10^{-19}.
	\]
	
	Константа образования комплекса:
	\[
	\beta = 10^{13}.
	\]
	
	\quad Большое значение $\beta$ означает высокую устойчивость комплекса [Cu(NH$_3$)$_4$]$^{2+}$, что приводит к полному растворению осадка Cu(OH)$_2$ в избытке аммиака.
	
	\textbf{Вывод.}  
	Образуется устойчивый комплекс тетраамминмеди(II).
	
	\subsubsection*{Опыт 9. Образование и разрушение комплексного аниона}
	
	\quad\textbf{ Ход опыта.}  
	К раствору Hg(NO$_3$)$_2$ добавляют KI → выпадает HgI$_2$.  
	Добавляют избыток KI → образуется растворимый комплекс.  
	Затем добавляют Na$_2$S → выпадает HgS.
	
	\textbf{Наблюдения.}  
	Осадок растворяется, затем снова выпадает.
	
	\textbf{Уравнения:}
	\[
	Hg^{2+} + 2I^- \rightarrow HgI_2\downarrow
	\]
	\[
	HgI_2 + 2I^- \rightarrow [HgI_4]^{2-}
	\]
	\[
	[HgI_4]^{2-} + S^{2-} \rightarrow HgS\downarrow + 4I^-
	\]
	
	\textbf{Расчёт равновесий.}
	
	Произведения растворимости:
	\[
	\text{ПР}_{\text{HgI}_2} = 10^{-29}, \qquad
	\text{ПР}_{[\text{HgI}_4]^{2-}} = 10^{-30}, \qquad
	\text{ПР}_{\text{HgS}} = 10^{-45}.
	\]
	
	Поскольку:
	\[
	10^{-30} \ll 10^{-29},
	\]
	HgI$_2$ растворяется в избытке I$^-$ благодаря образованию комплекса.
	
	И далее:
	\[
	10^{-45} \ll 10^{-30},
	\]
	образование HgS является наиболее выгодным по равновесию.
	
	\quad HgI$_2$ растворяется в избытке KI, но добавление S$^{2-}$ приводит к образованию чрезвычайно малорастворимого HgS.
	
	\textbf{Вывод.}  
	Избыток лиганда переводит осадок в раствор, но добавление S$^{2-}$ разрушает комплекс, образуя малорастворимый HgS.
	
\end{document}

